% --- Packages and PDF configuration ---

\usepackage[a5paper,left=12mm,right=12mm,top=12mm,bottom=6mm,includefoot,footskip=8mm]{geometry}
\usepackage{fontspec}
\usepackage{babel}
\babelprovide[import]{russian}
\usepackage{graphicx}
\usepackage{xcolor}
\usepackage{array}
\usepackage{longtable}
\usepackage{colortbl}
\usepackage{fancyhdr}
\usepackage{hyperref}
\hypersetup{
  unicode=true,
  pdfpagelayout=TwoColumnRight,
  colorlinks=false,
  pdfborder={0 0 0}
}

% PDF outline hierarchy matches the unnumbered document sections.
\makeatletter
\def\toclevel@section{0}
\def\toclevel@subsection{1}
\makeatother

% --- Common configuration and styles ---

% Fonts
% Golos Text has no italic face. Use an explicit oblique for emphasis.
\setmainfont{golostext-400.ttf}[Path=fonts/,Ligatures=NoCommon,
  ItalicFont=golostext-400.ttf,ItalicFeatures={FakeSlant}]
\newfontfamily\GolosText{golostext-400.ttf}[Path=fonts/,Ligatures=NoCommon,
  ItalicFont=golostext-400.ttf,ItalicFeatures={FakeSlant}]
\newfontfamily\GolosTextMedium{golostext-500.ttf}[Path=fonts/,Ligatures=NoCommon]
\newfontfamily\JetBrainsMono{jetbrainsmono-400.ttf}[Path=fonts/,Ligatures=NoCommon]
\newfontfamily\TablerDice{tabler-300.ttf}[Path=fonts/]
\urlstyle{same}

% Colors
\definecolor{rulegray}{gray}{0.60}

% Page style
\pagestyle{fancy}
\fancyhf{}
\renewcommand{\headrulewidth}{0pt}
\renewcommand{\footrulewidth}{0pt}
\fancyfoot[C]{{\GolosText\fontsize{9}{11.2}\selectfont\thepage}}

% Paragraph style
\setlength{\parindent}{0pt}
\newlength{\ParagraphSpacing}
\setlength{\ParagraphSpacing}{7.5pt}
\setlength{\parskip}{\ParagraphSpacing}
\newcommand{\ParagraphStyle}{%
  \GolosText\fontsize{9}{12}\selectfont\raggedright
  \setlength{\parindent}{0pt}%
  \setlength{\parskip}{\ParagraphSpacing}}

% Heading spacing
\newlength{\SectionBeforeSpace}
\newlength{\SubsectionBeforeSpace}
\newlength{\SectionAfterSpace}
\newlength{\SubsectionAfterSpace}
\setlength{\SectionBeforeSpace}{12pt}
\setlength{\SubsectionBeforeSpace}{8.5pt}
\setlength{\SectionAfterSpace}{10.5pt}
\setlength{\SubsectionAfterSpace}{\ParagraphSpacing}

% Heading formatting
\newcommand{\HeadingStyle}{\GolosTextMedium\raggedright}
\setcounter{secnumdepth}{0}
\makeatletter
\newcommand{\StartHeading}[6]{%
  % For consecutive headings, use this heading's before-space instead of
  % the previous heading's after-space, which standard sectioning retains.
  \if@nobreak
    \vskip-\lastskip
    \vskip-\dimexpr#4\relax
  \fi
  \@startsection{#1}{#2}{#3}{#4}{#5}{#6}}
\newcommand{\SectionHeading}{%
  \StartHeading{section}{1}{0pt}%
    {-\SectionBeforeSpace}{\SectionAfterSpace}%
    {\HeadingStyle\fontsize{16}{19}\selectfont}}
% Major sections start on a new page. Starred sections do not force a page break.
\renewcommand{\section}{%
  \@ifstar{\SectionHeading*}{\clearpage\SectionHeading}}
\renewcommand{\subsection}{%
  \StartHeading{subsection}{2}{0pt}%
    {-\SubsectionBeforeSpace}{\SubsectionAfterSpace}%
    {\HeadingStyle\fontsize{12}{15}\selectfont}}
\makeatother

% --- Cover ---

\newcommand{\CoverTitle}[1]{%
  {\GolosTextMedium\fontsize{33}{42}\selectfont #1\par}}

\newcommand{\CoverSubtitle}[1]{%
  \vspace{28.5pt}
  {\GolosText\fontsize{13}{16}\selectfont #1\par}}

\newcommand{\CoverEditionLine}[1]{%
  \vfill
  {\GolosText\fontsize{9}{12}\selectfont #1\par}}

\NewDocumentEnvironment{Cover}{+b}{%
  \thispagestyle{empty}
  \AddToHookNext{shipout/foreground}{%
    \put(12mm,-12mm){%
      \begingroup
      \setlength{\fboxsep}{0pt}%
      \setlength{\fboxrule}{0.5pt}%
      \raisebox{-\height}[0pt][0pt]{\color{black}\fbox{%
        \begin{minipage}[t][\dimexpr\paperheight-24mm-2\fboxrule\relax][t]{\dimexpr\paperwidth-24mm-2\fboxrule\relax}
          \centering
          \vspace*{72mm}
          #1%
          \vspace*{10mm}
        \end{minipage}}}%
      \endgroup}}%
  \null
}{}

% --- Publication details ---

\newlength{\ImprintBlockSpace}
\setlength{\ImprintBlockSpace}{\dimexpr\ParagraphSpacing+3pt\relax}

\newenvironment{Imprint}{%
  \thispagestyle{empty}%
  \ParagraphStyle
  \ParagraphAlignment
}{\par}

% Dedicated imprint typography uses paragraph spacing and its own block spacing.
% The following paragraph supplies the gap after titles and labels exactly once.
\newcommand{\ImprintTitle}[1]{%
  {\GolosTextMedium\fontsize{12}{15}\selectfont #1\par}\nobreak}
\newcommand{\ImprintLabel}[1]{%
  \par\addvspace{\ImprintBlockSpace}%
  {\GolosTextMedium #1\par}\nobreak}

\newenvironment{ImprintBlock}{%
  \par\addvspace{\ImprintBlockSpace}%
}{\par}

\newcommand{\ImprintText}[1]{#1\par}

% --- Contents ---

% Entry spacing
\newcommand{\ContentsSubentryIndent}{1.5em}
\newcommand{\ContentsEntryGap}{0.5em}
\newcommand{\ContentsLeaderStep}{0.55em}
\setcounter{tocdepth}{2}

% Entry formatting
\newcommand{\ContentsEntry}[2]{%
  \noindent#1\nobreak\hspace{\ContentsEntryGap}\color{rulegray}%
  \leaders\hbox to \ContentsLeaderStep{.\hss}\hfill\hspace{\ContentsEntryGap}\color{black}#2\par}
\newcommand{\ContentsSubentry}[2]{%
  \ContentsEntry{\hspace*{\ContentsSubentryIndent}#1}{#2}}

% Automatic contents integration
\makeatletter
\renewcommand{\l@section}[2]{\ContentsEntry{#1}{#2}}
\renewcommand{\l@subsection}[2]{\ContentsSubentry{#1}{#2}}

\newcommand{\ContentsEntries}{%
  \begingroup
  \ParagraphStyle
  \@starttoc{toc}%
  \endgroup}
\makeatother

% --- Dictionary description ---

% Paragraph alignment and hyphenation
\newcommand{\ParagraphAlignment}{%
  \selectlanguage{russian}%
  \rightskip=0pt plus 2em
  \parfillskip=0pt plus 1fil
  \spaceskip=\fontdimen2\font
  \xspaceskip=0pt
  \tolerance=1000
  \doublehyphendemerits=5000
  \finalhyphendemerits=5000
  \lefthyphenmin=3
  \righthyphenmin=3\relax}

\newcommand{\TextParagraph}[1]{%
  \begingroup
  \ParagraphStyle
  \ParagraphAlignment
  % Keep each explanatory paragraph on one page without joining adjacent blocks.
  \interlinepenalty=10000\relax
  #1\par
  \endgroup}

% Place reference link targets above the first line so it remains visible after navigation.
\makeatletter
\newcommand{\ReferenceTarget}[1]{\Hy@raisedlink{\hypertarget{#1}{}}}
\makeatother

% Lists
\newlength{\ListIndent}
\newcommand{\NumberedLabelSep}{0.25em}

\makeatletter
\newenvironment{List}{%
  \ParagraphStyle
  \ParagraphAlignment
  % Standard lists restore rightskip from this register during setup.
  \@rightskip=\rightskip
  % Keep lines and items together, leaving the list boundaries breakable.
  \interlinepenalty=10000\relax
  \@itempenalty=10000\relax
  % After a heading, \@nbitem merges its gap with the heading's existing skip.
  % Preserve the paragraph gap separately, as for a following TextParagraph.
  \if@nobreak \vskip\parskip \fi
  \setlength{\ListIndent}{2.5em}%
  \begin{list}{}{%
    % The surrounding paragraph spacing already supplies the external gap.
    \setlength{\topsep}{0pt}%
    \setlength{\partopsep}{0pt}%
    \setlength{\itemsep}{0pt}%
    \setlength{\parsep}{0pt}%
    \setlength{\leftmargin}{\ListIndent}%
    \setlength{\rightmargin}{0pt}%
    \setlength{\itemindent}{0pt}%
    \setlength{\listparindent}{0pt}%
    \setlength{\labelwidth}{\dimexpr\ListIndent-\NumberedLabelSep\relax}}%
}{\end{list}}
\makeatother

\newcommand{\NumberedItem}[2]{%
  \setlength{\labelsep}{\NumberedLabelSep}%
  \item[#1.]#2\par}
\newcommand{\BulletItem}[1]{%
  \setlength{\labelsep}{0.4em}%
  \item[\textbullet]#1\par}

% Standard block spacing uses one shared topsep above and below the table.
\newlength{\TransliterationBlockGap}
\newlength{\TransliterationColumnGap}
\setlength{\TransliterationBlockGap}{7mm}
\setlength{\TransliterationColumnGap}{4mm}

\newenvironment{TransliterationTable}{%
  \ParagraphStyle
  \setlength{\topsep}{\ParagraphSpacing}%
  \setlength{\partopsep}{0pt}%
  \begin{trivlist}%
  \item\relax
  \renewcommand{\arraystretch}{1.2}%
  \noindent\ignorespaces
}{\unskip\end{trivlist}}

\newcommand{\TextTableHeading}[1]{{\GolosTextMedium #1}}

\newenvironment{TransliterationBlock}{%
  \begin{tabular}[t]{@{}l@{\hspace{\TransliterationColumnGap}}l@{}}%
    \TextTableHeading{Буква} & \TextTableHeading{Запись}\\[\TableHeaderAfterSpace]\hline
}{\end{tabular}}

\newcommand{\TransliterationRow}[2]{#1 & #2\\\hline}

% --- Dictionary table ---

% Dimensions and alignment parameters
\setlength{\arrayrulewidth}{0.3pt}
\arrayrulecolor{rulegray}
\newlength{\TableHeaderAfterSpace}
\setlength{\TableHeaderAfterSpace}{0.2mm}
\newlength{\TableColumnPadding}
\setlength{\TableColumnPadding}{1.15mm}
\newlength{\TableRowHeight}
\newlength{\TableRowDepth}
\newlength{\TableDiceVerticalGap}
\newlength{\TableCapitalHeight}
\newlength{\TableTextLift}
\newlength{\TableDiceLeftShift}
\setlength{\TableRowHeight}{5.87mm}
\setlength{\TableRowDepth}{1.5mm}

% Align the visible left contour of the die frame with the left contour of «К».
% The value includes the two fonts' measured side bearings at their output sizes.
\setlength{\TableDiceLeftShift}{-0.198524mm}
\newsavebox{\TableCapitalReferenceBox}
\newsavebox{\TableDiceReferenceBox}

% Reference measurements and vertical alignment
\newcommand{\TableInitializeVerticalAlignment}{%
  \sbox{\TableCapitalReferenceBox}{{\GolosText\fontsize{9}{11.2}\selectfont H}}%
  \setlength{\TableCapitalHeight}{\ht\TableCapitalReferenceBox}%
  \setlength{\TableDiceVerticalGap}{\dimexpr\TableRowHeight/2-2.25mm\relax}%
  \setlength{\TableTextLift}{\dimexpr\TableRowHeight/2-\TableRowDepth-\TableCapitalHeight/2\relax}}

% Dice rendering
\newcommand{\TableDiceSymbol}[1]{%
  \ifnum#1<1
    \PackageError{dictionary}{Dice value #1 is below 1}{}
  \else
    \ifnum#1>6
      \PackageError{dictionary}{Dice value #1 is above 6}{}
    \else
      \char\numexpr"F08A+#1\relax
    \fi
  \fi}

\newcommand{\TableDiceGlyph}[1]{%
  \begingroup
  \sbox{\TableDiceReferenceBox}{\resizebox*{!}{4.5mm}{%
    \TablerDice
    \TableDiceSymbol{#1}}}%
  \raisebox{\dimexpr-\TableRowDepth+\TableDiceVerticalGap+\dp\TableDiceReferenceBox\relax}[0pt][0pt]{\usebox{\TableDiceReferenceBox}}%
  \endgroup}

\newcommand{\TableDiceGlyphBox}[1]{\makebox[5mm][l]{\hspace*{\TableDiceLeftShift}\TableDiceGlyph{#1}}}

\newcommand{\TableDiceHeader}[1]{%
  \TableHeaderFont\hspace*{0.04mm}#1}

% Search and copy support in PDF
\newcommand{\TableDiceUnicodeHex}[1]{%
  \ifcase#1\or 2680\or 2681\or 2682\or 2683\or 2684\or 2685\fi}

\newcommand{\TableDiceActualTextStart}[4]{%
  \begingroup
  \edef\TableDiceActualTextCommand{\endgroup
    \noexpand\pdfextension literal direct
    {/Span << /ActualText <FEFF\TableDiceUnicodeHex{#1}\TableDiceUnicodeHex{#2}\TableDiceUnicodeHex{#3}\TableDiceUnicodeHex{#4}> >> BDC}}%
  \TableDiceActualTextCommand}

\newcommand{\TableDiceSearchText}[1]{%
  \makebox[0pt][l]{%
    \pdfextension literal direct {3 Tr}%
    {\JetBrainsMono\fontsize{9}{11.2}\selectfont #1}%
    \pdfextension literal direct {0 Tr}}}

\newcommand{\TableDice}[1]{\expandafter\TableDiceFour#1\TableDiceEnd}
\def\TableDiceFour#1#2#3#4\TableDiceEnd{%
  \TableDiceSearchText{#1#2#3#4}%
  \TableDiceActualTextStart{#1}{#2}{#3}{#4}%
  \TableDiceGlyphBox{#1}\TableDiceGlyphBox{#2}\TableDiceGlyphBox{#3}\TableDiceGlyphBox{#4}%
  \pdfextension literal direct {EMC}}

% Cells, rows and header
\newcommand{\TableBodyFont}{\GolosText\fontsize{9}{11.2}\selectfont}
\newcommand{\TableHeaderFont}{\GolosText\fontsize{8}{10}\selectfont}
\newcommand{\TableMonospaceFont}{\JetBrainsMono\fontsize{9}{11.2}\selectfont}
\newcommand{\TableRowStrut}{\rule[-\TableRowDepth]{0pt}{\TableRowHeight}}
\newcommand{\TableTextCell}[1]{\raisebox{\TableTextLift}[0pt][0pt]{{\TableBodyFont\mbox{#1}}}}
\newcommand{\TableMonospaceCell}[1]{\raisebox{\TableTextLift}[0pt][0pt]{{\TableMonospaceFont\mbox{#1}}}}

\newcommand{\TableDataRow}[6][]{%
  \TableRowStrut\mbox{\raisebox{\dimexpr\TableRowHeight-\TableRowDepth\relax}[0pt][0pt]{#1}\TableDice{#2}} &
  \TableTextCell{#3} &
  \TableMonospaceCell{#4} &
  \TableMonospaceCell{#5} &
  \TableMonospaceCell{#6}\\\hline}

\newcommand{\TableHeader}[5]{%
  \TableDiceHeader{#1} &
  {\TableHeaderFont #2} &
  {\TableHeaderFont #3} &
  {\TableHeaderFont #4} &
  {\TableHeaderFont #5}\\[\TableHeaderAfterSpace]\hline
  \endhead}

% Table environment and spread positioning
\newcommand{\TableStartOnLeftPage}{%
  \clearpage
  \ifodd\value{page}
    \thispagestyle{empty}\null\clearpage
  \fi}

\newenvironment{DictionaryTable}{%
  \setlength{\tabcolsep}{\TableColumnPadding}%
  \setlength{\LTleft}{0pt}%
  \setlength{\LTright}{0pt}%
  \setlength{\LTpre}{0pt}%
  \setlength{\LTpost}{0pt}%
  \begin{longtable}{@{\extracolsep{\fill}}p{25mm}p{30mm}p{30mm}p{16.5mm}p{10mm}@{}}
}{\end{longtable}}
